\section{OpenCode cells}\label{app:opencode-bench}

\sloppy

This appendix documents the six OpenCode cells in the matrix
(one per benchmark $\times$ two models, across all three benchmarks).
It records their motivation, the apparatus that drives them, the
six-cell results, and the three cross-tool questions they answer.

\subsection{Why a fourth tool}

The threats-to-validity disclosure in
Section~\ref{sec:threats} (\emph{``OpenSpec and Ralph were measured under
Cline, not their canonical drivers''}) names a runtime confound: the
\texttt{openspec} and \texttt{ralph} cells in the matrix are run
with Cline as the agent runtime, not the canonical drivers
(Cursor / Claude Code / OpenCode for OpenSpec; an external
agent runtime per iteration for Ralph). The token-cost cells, the
operator-experience scores, and the MCP-call counts therefore partly
capture Cline's framing rather than the upstream tool's own surface.

The OpenCode cells are the explicit counter-evidence. They
use OpenCode's own native headless runtime
(\texttt{opencode run}) with the same prompt fingerprints, the same
MCP servers, and the same artefact schema as the rest of the study,
so the reader can compare a canonical-pairing data point against the
Cline-pairing data points elsewhere in the table. Concretely:

\begin{itemize}
\item OpenCode is a single binary that ships its own LLM
  client, MCP picker, session store, and headless invocation mode
  (invoked as \texttt{opencode run -{}-model <provider/model>
  -{}-format json}).
\item No IDE extension, no runtime wrapper. The agent talks directly
  to the configured provider with the operator-supplied API key.
\item MCP servers are wired in a per-bench \texttt{opencode.json}
  (GitHub MCP for the single-file bench; GitHub plus Atlassian for the
  video-support bench; GitHub MCP for the blog bench, where it
  is also the incremental-commit channel), so the picker behaviour is
  comparable to Cline's but the framing overhead is absent.
\end{itemize}

\subsection{Cells and apparatus}

Table~\ref{tab:opencode-cells} lists the six OpenCode cells and the
helper script that drives each.

\begin{table}[h]
\centering
\caption{The six OpenCode cells and the helper script that drives
each, with observed wall times.}
\label{tab:opencode-cells}
\footnotesize
\begin{tabular}{@{}llllr@{}}
\toprule
Cell & Bench & Model & Helper script & Wall (plan + exec) \\
\midrule
7  & github-bootstrap       & Claude Opus 4.8   & \texttt{run\_cell7.sh opus}  & 98\,s + 235\,s = 333\,s \\
7  & github-bootstrap       & GPT-5.5  & \texttt{run\_cell7.sh gpt}   & 119\,s + 184\,s = 303\,s \\
8  & worker-video-feature   & Claude Opus 4.8   & \texttt{run\_cell8.sh opus}  & 293\,s + 350\,s = 643\,s \\
8  & worker-video-feature   & GPT-5.5  & \texttt{run\_cell8.sh gpt}   & 188\,s + 253\,s = 441\,s \\
12 & personal-blog-scaffold & Claude Opus 4.8   & \texttt{run\_cell12.sh opus} & 118\,s + 1724\,s = 1842\,s \\
12 & personal-blog-scaffold & GPT-5.5  & \texttt{run\_cell12.sh gpt}  & 123\,s + 1245\,s = 1368\,s \\
\bottomrule
\end{tabular}
\end{table}

The helper scripts implement the same six-step life-cycle as
\texttt{tools/run\_cell\{1..6\}.sh}: prep, plan-run, capture-plan,
exec-run, capture-exec, drop. Two things differ. First, the
\texttt{plan-run} and \texttt{exec-run} sub-commands invoke
\texttt{opencode run} directly rather than handing off to an IDE
extension; the operator is unattended for the full duration of each
leg. Second, the artefact set adds two OpenCode-specific files per
leg: \texttt{transcripts/opencode-stdout.log} (the streaming JSON
event log from \texttt{opencode run} with the JSON format flag) and
\texttt{transcripts/session-export.json} (the deterministic export
fetched via \texttt{opencode export} on the captured session id).
Token totals and per-tool-call cost figures come from
\texttt{session-export.json} rather than from a heuristic byte-volume
estimate, removing the Parameter~9 caveat that applies to the
Cline-driven cells. The full operator runbook (provider auth, MCP
auth, model-string verification, drop / redact / guard recipe) is at
\texttt{HANDOFF-OPENCODE-RUN.md} in the study root. The MCP picker
configurations are at \texttt{tools/opencode/opencode.simple.json}
and \texttt{opencode.complex.json}; the headless wrapper is at
\texttt{tools/opencode/opencode\_run.sh}.

\subsection{Per-cell weighted totals and per-parameter scores}

Table~\ref{tab:opencode-totals} gives the per-cell weighted totals and
Table~\ref{tab:opencode-per-param} the full per-parameter scores for
the six OpenCode cells.

\begin{table}[h]
\centering
\caption{Per-cell weighted totals (max 5.0) for the six
OpenCode cells. Across all three benchmarks OpenCode ranks third on
the cross-model average ($3.701$, vs OpenSpec's $3.798$ and Ralph's
$3.728$ in Table~\ref{tab:aggregate}): the blog bench has no Jira
ticket and no rich external MCP for OpenCode to use (Param~2 and
Param~6 sit at the no-external-source baseline), so its blog cell is
third under both models ($3.540$ cross-model). OpenCode wins the
video-support bench outright ($3.955$ cross-model vs Ralph's $3.855$
and OpenSpec's $3.830$).}
\label{tab:opencode-totals}
\begin{tabular}{lrrr}
\toprule
Bench & Opus~4.8 & GPT-5.5 & Cross-model avg \\
\midrule
\texttt{github-repo-bootstrap}   & 3.540 & 3.675 & 3.608 \\
\texttt{worker-video-feature}    & 3.925 & 3.985 & \textbf{3.955} \\
\texttt{personal-blog-scaffold}  & 3.505 & 3.575 & 3.540 \\
\midrule
three-benchmark average          & 3.657 & 3.745 & 3.701 \\
\bottomrule
\end{tabular}
\end{table}

\begin{table}[h]
\centering
\caption{Per-parameter scores for the six OpenCode cells
(single-file, video-support, and blog benches under both models).
Column abbreviations: ``s'' = single-file, ``c'' = video-support, ``b'' = blog;
``/O'' = Opus~4.8, ``/G'' = GPT-5.5. Anchor-5 highlights:
OpenCode anchors at $5$ on Task granularity (3) and Hand-off friction
(10) across all three benches and both models, on Test-plan coverage
(4) on all but the Opus blog cell, and on Code-reference accuracy (5)
and MCP integration (6) on the video-support bench; the blog bench
drops Param~2 (no external ticket) and Param~6 (GitHub-MCP-only) to
the no-external-source baseline, which is why the blog columns sit
below the video-support columns.}
\label{tab:opencode-per-param}
\setlength{\tabcolsep}{4pt}
\begin{tabular}{rp{5.2cm}cccccc}
\toprule
\# & Parameter & s/O & s/G & c/O & c/G & b/O & b/G \\
\midrule
1  & Artifact completeness                    & 2 & 2 & 2 & 2 & 2 & 2 \\
2  & Requirement traceability                 & 2 & 2 & 4 & 4 & 2 & 2 \\
3  & Task granularity                         & 5 & 5 & 5 & 5 & 5 & 5 \\
4  & Test-plan coverage                       & 5 & 5 & 5 & 5 & 4 & 5 \\
5  & Code-reference accuracy                  & 4 & 5 & 5 & 5 & 5 & 4 \\
6  & MCP / external-source integration        & 3 & 3 & 5 & 5 & 3 & 3 \\
7  & Reproducibility (artefact persistence)   & 4 & 4 & 4 & 4 & 4 & 4 \\
8  & Latency to ready-to-implement            & 5 & 5 & 5 & 5 & 5 & 5 \\
9  & Token / API cost                         & 2 & 4 & 1 & 3 & 2 & 3 \\
10 & Hand-off friction to implementation      & 5 & 5 & 5 & 5 & 5 & 5 \\
11 & Tool licensing \& accessibility          & 5 & 5 & 5 & 5 & 5 & 5 \\
12 & Learning curve / time-to-productivity    & 3 & 3 & 3 & 3 & 3 & 3 \\
13 & Human-in-the-loop validation experience  & 2 & 2 & 2 & 2 & 2 & 2 \\
14 & Autonomous capacity                      & 2 & 2 & 2 & 2 & 2 & 2 \\
15 & Runtime integration \& ergonomics        & 3 & 3 & 3 & 3 & 3 & 3 \\
16 & Code-review / hand-off review surface    & 4 & 4 & 4 & 4 & 4 & 4 \\
\midrule
   & \textbf{Weighted total}                  & \textbf{3.540} & \textbf{3.675} & \textbf{3.925} & \textbf{3.985} & \textbf{3.505} & \textbf{3.575} \\
\bottomrule
\end{tabular}
\end{table}

\subsection{Three cross-tool questions, answered}

\paragraph{Q1.\ Does the Cline confound move the rubric ranking?}
\textbf{No --- where the canonical and Cline pairings are directly
comparable (the externally-grounded video-support bench) the native pairing
\emph{out-scores} the Cline-driven tools; the cross-tool average
leader changes only because of benchmark composition, not the driver.}
Compare OpenCode's per-cell weighted total against the
\texttt{openspec} and \texttt{ralph} totals in
Table~\ref{tab:aggregate}: on the video-support bench OpenCode leads under
both models (Opus~4.8 $3.925$ vs \texttt{openspec} $3.870$ vs
\texttt{ralph} $3.895$; GPT-5.5 $3.985$ vs $3.790$ vs $3.815$). On the
single-file bench under Opus~4.8 OpenCode ($3.540$) sits \emph{below}
\texttt{openspec} ($3.695$); on the blog bench OpenCode is
third under both models ($3.505$/$3.575$) because the empty repository
strips the external grounding it depends on. The cross-model
three-benchmark average therefore favours OpenSpec ($3.798$) over
Ralph ($3.728$) and OpenCode ($3.701$) --- but this is a
benchmark-composition effect: on the one bench where integration
demand makes the canonical vs Cline pairing directly comparable,
OpenCode's native pairing is on top, so the Cline-pairing confound is
not \emph{inflating} OpenSpec's or Ralph's scores past OpenCode's. The
Cline-pairing caveats (Cline's prompt-assembly overhead inflates
Param~9 cost estimates; Cline's chat surface partly drives the OPX
scores; Cline's MCP auto-injection partly drives the Param~6 MCP-call
counts) remain literally true, but the cross-tool average leader is
determined by the benchmark composition --- in particular the
breadth-oriented blog bench --- not by the driver.

\paragraph{Q2.\ Does the canonical pairing change Parameter~9 (Token
/ API cost)?} \textbf{Yes --- and the change is consistent with the
``rate-card overestimates real spend on heavy-cache workloads''
hypothesis.} The OpenCode session export under Opus reports billed
tokens deterministically, and the GPT-5.5 cells are reconstructed by
the per-line-rate backfill in \texttt{opencode\_run.sh} so each token
line item (input, output, cached read, cached write) is priced
separately. OpenCode's cache-aware totals are materially
\emph{lower} than the rate-card estimates would predict despite
OpenCode emitting the most tokens of any tool: video-support bench Opus~4.8
is \$22.07 cache-aware vs rate-card-equivalent \$286.94, and the
blog bench Opus~4.8 is \$49.57 vs a rate-card-equivalent
$\sim$\$763 (its 40.1\,M tokens at \$19/M). The delta is the framing
overhead \emph{plus} the prompt-caching discount that the rate-card
path does not model. A second reading specific to the Opus~4.8 column:
although the cache-aware figures are far below rate-card, their
\emph{absolute} level is the highest in the study --- the blog
cell at \$49.57 (vs \$8.38 for the same cell under GPT-5.5) is the
single most expensive cell in the corpus, because the Opus~4.8
OpenCode greenfield run is a 40.1\,M-token deep-context session.
OpenCode is thus the cost outlier under Opus~4.8 on every bench and
the cheapest tool on the video-support bench under GPT-5.5 --- the cost
ordering across tools cannot be read without specifying the model.

\paragraph{Q3.\ Does OpenCode reproduce the regex-passthrough
trap on the video-support bench?} \textbf{Only under GPT-5.5 --- under
Opus~4.8 OpenCode \emph{avoids} it.} The trap (failure to extend
\texttt{STANDARD\_PATH\_REGEX} from \texttt{[a-z]+} to
\texttt{[a-z0-9]+} so that \texttt{mp4|webm|mov} would route
correctly) is the one cell in the matrix where frontier-model
substitution changes the outcome. Under \textbf{Opus~4.8} the
OpenCode planning leg diagnoses the defect directly --- the plan body
states that ``\texttt{mp4} contains a digit and fails the regex
\ldots\ the group must become \texttt{([a-z0-9]+)}''
(\texttt{OPENCODE\_PLAN.md}~\S 2.3) and Task~2 broadens the character
class before any code is written, so the trap never fires. Under
\textbf{GPT-5.5} the same tool asserts the regex was already permissive
(plan body: \texttt{[a-z]+} ``already allows mp4/webm/mov''), falls in,
and the execution leg recovers via a Vitest failure during apply by
widening the character class one character. Recovery imposed no test
regressions. The defect-convergence story in
Section~\ref{subsec:defect-convergence} therefore reads: \emph{the
spec format avoids the trap under both models without inspecting the
digit (OpenSpec), while plan-content quality can avoid it only when it
is high enough --- and is model-dependent}. OpenCode's Opus~4.8 plan
is the single most precise code-reference diagnosis in the study and
avoids the trap, but the same tool falls in under GPT-5.5; only the
4-artefact spec scaffold front-loads ``parser broadening'' as a
first-class requirement under both models.

\subsection{Cross-model behavioural asymmetry on OpenCode}

A cross-model behavioural difference surfaces on the single-file bench
OpenCode cells: \emph{precondition-mismatch resolution scope}. The
single-file bench prompt asks the planner to verify
\texttt{terraform fmt -check -recursive} runs clean, and the target
repo carries 13 pre-existing unformatted files unrelated to the
requested change. OpenCode/Opus deliberately scoped its fix to the
new file (recording the global drift as a pre-existing condition in
\texttt{deviations.json}, exit-0 verifying only the target
subdirectory); OpenCode/GPT-5.5 ran a repo-wide
\texttt{terraform fmt -recursive} to clear all 13 files of drift
before re-running the strict-clean postcondition. Both responses
produce a green strict-clean postcondition; neither violates the
test plan. The same scope-resolution split surfaces on the OpenSpec
and Ralph single-file bench cells under GPT-5.5 (and on neither of
their Opus equivalents), confirming this is a model-bound asymmetry
that generalises across all canonical-pairing tools rather than a
tool-specific OpenCode behaviour: GPT-5.5 pulls scope outward
toward strict-postcondition closure, Opus pulls scope inward toward
single-file-change minimality. Recorded in
Appendix~\ref{app:procdev}.4 with the corresponding diff-stat
asymmetry in Table~\ref{tab:metadata-exec}.
